% Figure: a spool pulled by a string wound under its inner hub. The
% direction the spool rolls follows from the moment about the contact
% point; a pull whose line of action passes through the contact gives
% no rolling at all, the critical-angle case.
\begin{figure}[htbp]
\centering
\begin{tikzpicture}[scale=1.0,
  frc/.style={draw=engblue, -{Stealth}, very thick},
  lbl/.style={font=\scriptsize},
]
% left panel: horizontal pull under the hub, spool rolls toward the pull
\begin{scope}[shift={(1.2,1.3)}]
  \draw[engblack, thick] (-1.8,-1.3) -- (2.6,-1.3);
  \draw[engblue, thick, fill=engblue!5] (0,0) circle (1.3);   % outer radius R
  \draw[enggray, thick, fill=enggray!12] (0,0) circle (0.6);  % inner hub r
  \fill[engblack] (0,0) circle (1.2pt);
  \fill[engred] (0,-1.3) circle (1.6pt);
  \node[lbl, engred, anchor=north] at (0,-1.35) {$C$};
  % string leaves the bottom of the hub, pulled horizontally right
  \draw[frc] (0,-0.6) -- (2.3,-0.6) node[right, lbl, engblue] {$P$};
  % moment arm from C up to the line of action
  \draw[enggray, dashed] (0,-1.3) -- (0,-0.6);
  \node[lbl, enggray, anchor=east] at (-0.05,-0.95) {$R{-}r$};
  \node[lbl, engamber, anchor=south] at (0,1.45) {horizontal pull};
  \draw[-{Stealth}, engamber, thick] (-0.75,-1.62) arc (150:30:0.5);
  \node[lbl, engamber, anchor=north] at (0,-1.75) {rolls toward $P$};
\end{scope}
% right panel: pull aimed through the contact point, no rolling
\begin{scope}[shift={(6.2,1.3)}]
  \draw[engblack, thick] (-1.8,-1.3) -- (2.6,-1.3);
  \draw[engblue, thick, fill=engblue!5] (0,0) circle (1.3);
  \draw[enggray, thick, fill=enggray!12] (0,0) circle (0.6);
  \fill[engblack] (0,0) circle (1.2pt);
  \fill[engred] (0,-1.3) circle (1.6pt);
  \node[lbl, engred, anchor=north] at (0,-1.35) {$C$};
  % string leaves the bottom of the hub aimed so its line passes through C
  % hub-bottom at (0,-0.6); aim the line through C at (0,-1.3): direction
  % points down-and-right toward a far point; draw pull along that line
  \draw[frc] (0,-0.6) -- (2.03,-1.90) node[right, lbl, engblue] {$P$};
  % show the line continued to C
  \draw[enggray, dashed] (0,-0.6) -- (0,-1.3);
  \node[lbl, engamber, anchor=south] at (0,1.45) {pull aimed through $C$};
  \node[lbl, engamber, anchor=north] at (0,-1.75) {no rolling};
\end{scope}
\end{tikzpicture}
\caption[A spool pulled by a string wound on its inner hub]{A spool of
outer radius $R$ pulled by a string wound on an inner hub of radius $r$.
The direction of rolling follows from the moment of the pull about the
contact point $C$, taken as the instantaneous centre. A horizontal pull
off the bottom of the hub (left) has a line of action a height $R-r$
above the contact, so its moment $P(R-r)$ rolls the spool toward the
pull, winding the string up as it comes, the result that surprises most
first-time solvers. Tilt the pull until its line passes through the
contact (right) and the moment arm vanishes: the spool neither rolls
forward nor back but is on the verge of sliding. The lesson is to take
moments about $C$ rather than about the centre, because the friction
force at $C$ then drops out of the moment equation and the sense of
rolling reads straight off the geometry of the pull.}
\label{fig:vol03ch04:pulled-spool}
\end{figure}
